\documentclass[12pt]{article}

% a template that a friend gave, it's worked well enough for me
% i have added some packages and stuff that have proved useful

\usepackage{fancyhdr}
\usepackage{tipa}
\usepackage{fontspec}
\usepackage{amsfonts}
\usepackage{enumitem}
\usepackage[margin=1in]{geometry}
\usepackage{graphicx}
\usepackage{float}
\usepackage{amsmath}
\usepackage{braket}
\usepackage{amssymb}
\usepackage{booktabs}
\usepackage{hyperref}
\usepackage{mathtools}
\usepackage{xcolor}
\usepackage{float}
\usepackage{algpseudocodex}
\usepackage{titlesec}
\usepackage{bbm}

\pagestyle{fancy}
\fancyhf{} % sets both header and footer to nothing
\lhead{Kevin Sheng}
\setmainfont{Comic Neue}
\renewcommand{\headrulewidth}{1pt}
\setlength{\headheight}{0.75in}
\setlength{\oddsidemargin}{0in}
\setlength{\evensidemargin}{0in}
\setlength{\voffset}{-.5in}
\setlength{\headsep}{10pt}
\setlength{\textwidth}{6.5in}
\setlength{\headwidth}{6.5in}
\setlength{\textheight}{8in}
\renewcommand{\headrulewidth}{0.5pt}
\renewcommand{\footrulewidth}{0.3pt}
\setlength{\textwidth}{6.5in}
\usepackage{setspace}
\usepackage{multicol}
\usepackage{float}
\setlength{\columnsep}{1cm}
\setlength\parindent{24pt}
\usepackage [english]{babel}
\usepackage [autostyle, english = american]{csquotes}
\MakeOuterQuote{"}

\setlength{\parskip}{6pt}
\setlength{\parindent}{0pt}

\titlespacing\section{0pt}{12pt plus 4pt minus 2pt}{0pt plus 2pt minus 2pt}
\titlespacing\subsection{0pt}{12pt plus 4pt minus 2pt}{0pt plus 2pt minus 2pt}
\titlespacing\subsubsection{0pt}{12pt plus 4pt minus 2pt}{0pt plus 2pt minus 2pt}

\hypersetup{colorlinks=true, urlcolor=blue}

\newcommand{\correction}[1]{\textcolor{red}{#1}}

\usepackage{nicematrix}

\rhead{Math 110A}

\makeatletter
\def\@seccntformat#1{%
  \expandafter\ifx\csname c@#1\endcsname\c@section\else
  \csname the#1\endcsname\quad
  \fi}
\makeatother

\newcommand{\lra}{\xLeftrightarrow}
\newcommand{\ra}{\xRightarrow}
\DeclareMathOperator{\pr}{pr}
\DeclareMathOperator{\Mat}{Mat}

\begin{document}

\section{Problem 1}

Fix any basis on $S$.
I propose that our $h$ should map any $s$ to the matrix associated with
the linear transformation $T_s(t)=s \cdot t$.
We can multiply two elements in $S$ since it's also a $\C$-algebra.

(that i can do this is just 115a stuff,
like take each of the basis vectors and see how they change to get the matrix)

Let's use $v: S \to \C^n$ to represent the unique linear combination
on $S$'s basis that adds up to its input.
By construction, $h(s)v(s')=v(s \cdot s')$.


Now we prove this is a homomorphism of algebras:
\begin{enumerate}
  \item \textbf{Addition:}
        \begin{align*}
          h(s_1+s_2)v(s)
           & = v((s_1+s_2)s)         \\
           & = v(s_1s+s_2s)          \\
           & = v(s_1s)+v(s_2s)       \\
           & = h(s_1)v(s)+h(s_2)v(s)
        \end{align*}
        Since this is true for all $s$, we could set it equal to each of the basis
        vectors to see that, say, the first column should be equal to the sum of
        the other two first columns, which is indeed the case.

  \item \textbf{Multiplication:}
        \begin{align*}
          h(s_1 \cdot s_2)v(s)
           & = v(s_1s_2s)    \\
           & = h(s_1)v(s_2s) \\
           & = h(s_1)h(s_2)s
        \end{align*}

  \item \textbf{One:}
        \[h(1)v(s)=v(s) \implies h(1)=I\]

  \item \textbf{Preservation:}
        \[h(c)v(s)=v(cs) \implies h(c)=cI\]
\end{enumerate}

With that shown, now we need that it's injective.
Let $s_1, s_2 \in S$ be s.t. $h(s_1)=h(s_2)$ and by extension $T_{s_1}=T_{s_2}$.
Given this $T_{s_1}(1)=T_{s_2}(1) \implies s_1=s_2$. $\square$

\pagebreak

\section{Problem 2}

Consider the following $\C$-subalgebras:
\begin{enumerate}
  \item \[
          \left(
          \left\{
          \begin{bmatrix}
            c & 0 \\
            0 & c
          \end{bmatrix}
          \mathrel{\Big|} c \in \mathbb{C}
          \right\},
          \quad
          c \mapsto
          \begin{bmatrix}
            c & 0 \\
            0 & c
          \end{bmatrix}
          \right)
        \]

  \item \[
          \left(
          \left\{
          \begin{bmatrix}
            a & b \\
            c & d
          \end{bmatrix}
          \mathrel{\Big|} a,b,c,d \in \mathbb{C}
          \right\},
          \quad
          c \mapsto
          \begin{bmatrix}
            c & 0 \\
            0 & c
          \end{bmatrix}
          \right)
        \]

  \item \[
          \left(
          \left\{
          \begin{bmatrix}
            a & 0 \\
            0 & b
          \end{bmatrix}
          \mathrel{\Big|} a,b \in \mathbb{C}
          \right\},
          \quad
          c \mapsto
          \begin{bmatrix}
            c & 0 \\
            0 & c
          \end{bmatrix}
          \right)
        \]

  \item \[
          \left(
          \left\{
          \begin{bmatrix}
            a & b \\
            0 & a
          \end{bmatrix}
          \mathrel{\Big|} a,b \in \mathbb{C}
          \right\},
          \quad
          c \mapsto
          \begin{bmatrix}
            c & 0 \\
            0 & c
          \end{bmatrix}
          \right)
        \]
        Multiplication in this set is closed:
        \[\begin{bmatrix}
            a & b \\
            0 & a
          \end{bmatrix}\begin{bmatrix}
            a' & b' \\
            0  & a'
          \end{bmatrix}
          = \begin{bmatrix}
            aa' & ab'+a'b \\
            0   & aa'
          \end{bmatrix}\]

  \item \[
          \left(
          \left\{
          \begin{bmatrix}
            a & b \\
            0 & c
          \end{bmatrix}
          \mathrel{\Big|} a,b,c \in \mathbb{C}
          \right\},
          \quad
          c \mapsto
          \begin{bmatrix}
            c & 0 \\
            0 & c
          \end{bmatrix}
          \right)
        \]
\end{enumerate}
I only verified closedness for the fourth one because it implies the closedness
of the fifth subalgebra and all the others are fairly obvious.

All pairs of subalgebras besides 3 and 4 have different dimensions as vector spaces.
Thus, it is impossible to create a bijection between the elements of these pairs
that is still a linear transformation.

The only case that remains is 3 and 4, since they both have dimension $2$.
Take any homomorphism $h$.
Let its images on the two spanning matrices be like so:
\begin{align*}
  h\left(\begin{bmatrix}1 & 0 \\ 0 & 0\end{bmatrix}\right)=\begin{bmatrix}x & y \\0 & x\end{bmatrix}
   &  & h\left(\begin{bmatrix}0 & 0 \\ 0 & 1\end{bmatrix}\right)=\begin{bmatrix}x' & y' \\0 & x'\end{bmatrix}
\end{align*}
Since $h$ is a homomorphism,
\[h\left(\begin{bmatrix}1 & 0 \\ 0 & 0\end{bmatrix}\right)
  +h\left(\begin{bmatrix}0 & 0 \\ 0 & 1\end{bmatrix}\right)
  =h(I_2)=I_2 \implies
  x+x'=1,y+y'=0\]
and we can write
\[h\left(\begin{bmatrix}0 & 0 \\ 0 & 1\end{bmatrix}\right)=\begin{bmatrix}1-x & -y \\0 & 1-x\end{bmatrix}\]

Now consider the multiplication of these two matrices:
\begin{align*}
             & h\left(\begin{bmatrix}1 & 0 \\ 0 & 0\end{bmatrix}\right)
  \cdot h\left(\begin{bmatrix}0 & 0 \\ 0 & 1\end{bmatrix}\right) = h(\mathbf{0}) = \mathbf{0}                     \\
  \implies{} & \begin{bmatrix}x & y \\0 & x\end{bmatrix}\begin{bmatrix}1-x & -y \\0 & 1-x\end{bmatrix}=\mathbf{0} \\
  \implies{} & \begin{bmatrix}x(1-x) & y-2xy \\ 0 & x(1-x)\end{bmatrix} = \mathbf{0}                              \\
  \implies{} & x \in \{0, 1\}, y=0
\end{align*}
Thus, for any matrix in the third subalgebra,
\begin{align*}
  h\left(\begin{bmatrix}a & 0 \\ 0 & b\end{bmatrix}\right)
   & = ah\left(\begin{bmatrix}1 & 0 \\ 0 & 0\end{bmatrix}\right)
  +bh\left(\begin{bmatrix}0 & 0 \\ 0 & 1\end{bmatrix}\right)     \\
   & = a\begin{bmatrix}x & 0 \\0 & x\end{bmatrix}+b\begin{bmatrix}1-x & -y \\0 & 1-x\end{bmatrix} \\
   &= \begin{bmatrix}ax+b(1-x) & 0 \\ 0 & ax+b(1-x)\end{bmatrix}
\end{align*}
As you can see, no matrix is mapped to another one with a nonzero top right element.
Thus, no homomorphism from the third to the fourth subalgebra can simultaneously
be a bijection.

\pagebreak

\section{Problem 3}

Let $\phi_1$ and $\phi_2$ be the associated homomorphisms for $S_1$ and $S_2$ respectively.

It suffices to show that for all $r \in R$ and $s \in S_1$ we have $f(s) \cdot \phi_2(r) \in f(S_1)$.

Note that since $f$ is a homomorphism of algebras, $\phi_2(r)=f(\phi_1(r))$ and
\[f(s) \cdot \phi_2(r)=f(s) \cdot f(\phi_1(r))=f(s \cdot \phi_1(r)) \in f(S_1)\quad\square\]

\section{Problem 4}

I'll first show that $T=\left\{\sum_{i \ge 0} r_i s^i\right\}$ is an $R$-subalgebra.

To do this, we show that it's closed under addition:
\[\sum_{i \ge 0} r_i s^i + \sum_{i \ge 0} r'_i s^i
  = \sum_{i \ge 0} (r_i+r'_i) s^i\]
multiplication:
\[\sum_{i \ge 0} r_i s^i \cdot \sum_{i \ge 0} r'_i s^i
  = \sum_{i \ge 0} \left(\sum_{j=0}^{i} r_i r'_{i-j}\right) \cdot s^i\]
and scaling:
\[r \cdot \sum_{i \ge 0} r_i s^i = \sum_{i \ge 0} (r \cdot r_i) \cdot s^i\]

Now consider any other $R$-algebra $T'$ that contains $s$.
It certainly must have at least $3$ elements: $0$, $1$, and $s$ itself.
Since it's closed under scaling, $\phi(r) \cdot 1=\phi(r)$ must be in $T$
for all $r$ as well.

$T'$ is a subring, so we should be able to add and multiply these elements in any
way and still stay in $T'$.
"Any way" here includes the linear combination-style things that we were doing,
so $\sum_{i \ge 0} r_i s^i \in T'$ and by extension $T \subseteq T'$. $\square$

\pagebreak

\section{Problem 5}

Let $D_m$ for $m \ge 0$ be defined like so:
\[(D_m)_{ij}=\begin{cases}
    1 & j=i+m            \\
    0 & \text{otherwise}
  \end{cases}\]
As one can see, $D_1=X$.

To get the subspace generated by $X$, I'll first prove that $(D_1)^m=D_m$.
We proceed by induction, and the base case for $k=1$ is near tautological.

If we assume the truth of the statement for $m-1$, then for $m < n$
\begin{align*}
  \left((D_1)^m\right)_{ij}
   & = \sum_{k=1}^{n} (D_1)^{m-1}_{ik} (D_1)_{kj}          \\
   & = \sum_{k=1}^{n} (D_{m-1})_{ik} (D_1)_{kj}            \\
   & = \sum_{k=1}^{n} \mathbf{1}[k=i+m-1]\mathbf{1}[j=k+1] \\
   & = \mathbf{1}[j=i+m]
\end{align*}
which is exactly what we wanted.

Notice that for $m \ge n$, $D_m=\mathbf{0}^{n \times n}$,
so we don't need to consider any powers at or above $n$.
Thus, the generated subspace can be expressed like so:
\[S=\left\{\sum_{i=0}^{n-1} r_i D_{i+1}\right\}\]
and any matrix in $S$ can be written like
\[M = \begin{bNiceMatrix}[columns-width=auto]
    r_0 & r_1 & r_2    & \cdots & r_{n-2} & r_{n-1} \\
        & r_0 & r_1    & r_2    & \ddots  & r_{n-2} \\
        &     & \ddots & \ddots & \ddots  & \vdots  \\
        &     &        & r_0    & r_1     & r_2     \\
        &     &        &        & r_0     & r_1     \\
        &     &        &        &         & r_0
  \end{bNiceMatrix}\]
Since we have $n$ linearly independent vectors, the dimension of $S$ is $n$.

\section{Problem 6}

I'll just show that the only homomorphisms that go from $\Q^I \to \Q$ are
the projection functions
\[\pr_j((q_i)_{i \in I})=q_j\]
To prove part A, we can set $I=\{1, 2\}$.

For the sake of convenience, I'll let $e_i$ denote the $I$-tuple
that has all of its entries set to $0$ except for the $i$th entry,
which is set to $1$.

Now consider any homomorphism $h: \Q^I \to \Q$.
We must have $h((0)_{i \in I})=0$ and $h((1)_{i \in I})=1$.

Notice that for any $e_i$
\begin{align*}
  e_i=e_i \cdot e_i
   & \implies h(e_i)=h(e_i \cdot e_i) \\
   & \implies h(e_i)=h(e_i)^2         \\
   & \implies h(e_i) \in \{0, 1\}
\end{align*}
We also have
\[h((1)_{i \in I})=\sum_{i \in I} h(e_i)=1\]
These two conditions force there to be a unique $\hat{i} \in I$ s.t.
\[h(e_i)=\begin{cases}
    1 & i=\hat{i}        \\
    0 & \text{otherwise}
  \end{cases}\]
If there weren't exactly one, then the summation equality above would break.

Now we can break down any mapping like so:
\begin{align*}
  h((q_i)_{i \in I})
   & = h\left(\sum_{i \in I} q_i \cdot e_i\right) \\
   & = \sum_{i \in I} q_i h(e_i)                  \\
   & = q_{\hat{i}} h(e_{\hat{i}})                 \\
   & = q_{\hat{i}}
\end{align*}
so any homomorphism $h: \Q^I \to \Q$ must be equal to some $\pr_i$. $\square$

\end{document}
